\documentclass[10pt,letterpaper]{article}

\usepackage[utf8]{inputenc}
\usepackage[T1]{fontenc}
\usepackage{lmodern}
\usepackage{amsmath,amssymb,amsfonts}
\usepackage{booktabs}
\usepackage{graphicx}
\usepackage{geometry}
\usepackage{hyperref}
\usepackage{xcolor}
\usepackage{array}
\usepackage{multirow}

\geometry{
  letterpaper,
  top=0.8in,
  bottom=0.85in,
  left=0.72in,
  right=0.72in,
  columnsep=0.28in
}

\hypersetup{
  colorlinks=true,
  linkcolor=blue!70!black,
  citecolor=blue!70!black,
  urlcolor=blue!70!black,
  pdftitle={Multimodal Narrative Systems Engineering: Signal-Grounded Measurement and an Exploratory State-Space Model for Audiovisual Media},
  pdfauthor={Jair Ivan Rojas Ruiz (ORCID: 0009-0003-1109-4583)}
}

\title{\textbf{Multimodal Narrative Systems Engineering: Signal-Grounded Measurement and an Exploratory State-Space Model for Audiovisual Media}}

\author{
  \textbf{Jair Ivan Rojas Ruiz} (\emph{Jair Rojas})\\
  Founder \& Principal Investigator, Narrative Systems Engineering (NSE)\\
  ORCID: \href{https://orcid.org/0009-0003-1109-4583}{\texttt{0009-0003-1109-4583}} \quad|\quad
  Concept DOI: \href{https://doi.org/10.5281/zenodo.23987655}{\texttt{10.5281/zenodo.23987655}}\\
  Repository: \href{https://github.com/jaiborojas-cloud/NarrativeSystemsEngineering}{\texttt{github.com/jaiborojas-cloud/NarrativeSystemsEngineering}}
}

\date{November 6, 2026 \quad(\texttt{PAPER-001} Preprint v1.1; supersedes v1.0 of November 5, 2026)\\\small Copyright 2026 Jair Ivan Rojas Ruiz. All rights reserved.}

\begin{document}

\maketitle

\begin{abstract}
\noindent
Narratives across film, television, political discourse, and digital media are typically analyzed either through qualitative frameworks or through text-only natural language processing, and scores produced by prompting large language models on arbitrary $0\text{--}100$ scales are hard to trace to the audiovisual signal. This preprint presents \textbf{Narrative Systems Engineering (NSE)} and its computational core, \textbf{Computational Narrative Dynamics (CND)}, an exploratory framework originated by Jair Ivan Rojas Ruiz \cite{RojasRuiz2026NSE} with two layers that we keep explicitly separate. The first layer is \emph{measurement}: acoustic level and speech activity from PCM audio, visual change between sampled frames (Jensen-Shannon divergence of color distributions), luminance, saturation and color entropy, and transcript-based speech rate and lexical descriptors, each linked to exact time spans and source hashes. The second layer is a \emph{model}: composite indices (Narrative Energy, Entropy, Complexity, Inertia) with author-chosen weights, a single-coder stance coding $z(t)$, and a Friedkin-Johnsen-style state-space model $x(t+1) = f(x(t), u(t))$ of how these inputs might move audience belief states. We apply the framework descriptively to three short works ($513.6\text{ s}$, $77$ windows): a 1994 Spanish university debate, the 1964 \emph{``Daisy''} spot, and the 1986 \emph{Challenger} address. The measurement layer yields reproducible, time-aligned profiles that differ visibly across the three works. The model layer is not validated: its weights are uncalibrated, the stance coding has no measured inter-annotator agreement, the counterfactual montage results are simulations of the model itself, and no audience response has been measured.
\end{abstract}

\noindent\textbf{Keywords:} Narrative Systems Engineering, Computational Narrative Dynamics, Multimodal Signal Processing, Shannon Narrative Entropy, Narrative Inertia, State-Space Model, Exploratory Model.

\section{Introduction}

Human societies are structured, stabilized, and transformed by the transmission of symbolic narratives across audiovisual and linguistic media \cite{Green2000,Jones2010,Shiller2017}. However, bridging the gap between raw media signals (pixels, acoustic waveforms, lexical tokens) and higher-order cognitive and social dynamics remains an open scientific and engineering problem.

Existing computational approaches divide into two disconnected paradigms:
\begin{enumerate}
  \item \textbf{Low-level media descriptors:} Computer vision and audio pipelines extract shot boundaries, color histograms, or automatic speech transcripts without modeling how these observables couple to form persuasive or symbolic structures \cite{Vicol2018}.
  \item \textbf{Uncalibrated generative summaries:} Prompting LLMs to output ``narrative energy'' or ``belief shift'' without formal measurement operators produces plausible-sounding numerical fiction. For instance, defining Narrative Energy as a naive product $E = \mathrm{Intensity} \times \mathrm{Coherence} \times \mathrm{Repetition} \times \mathrm{Identification}$ fails dimensional analysis and collapses to zero whenever any single factor vanishes.
\end{enumerate}

We adopt four design criteria for every proposed quantity; they are aims, not evidence of validity: (i) \emph{Signal Grounding} in physical PCM16 audio samples, RGB24 rasters, and time-aligned speech tokens; (ii) \emph{Stated Units and Bounds} (Shannon bits \cite{Shannon1948,Lin1991}, $\mathrm{dBFS}$, and bounded state spaces $[-1,1]^K$); (iii) \emph{Layer Separation} between observable work material, symbolic interpretation, and state response; and (iv) \emph{Executable Reproducibility}.

\textbf{Terminology note:} the phrase ``Narrative Systems Engineering'' is also used elsewhere (game-design notes, an organizational coaching framework, and Strangium's November 2, 2026 \texttt{NSE-001} screenplay taxonomy). This work uses it for signal-level measurement of audiovisual works combined with an exploratory state-space model. The founding documents are archived since September 8, 2026 and deposited in Zenodo \cite{RojasRuiz2026NSE}; these records document authorship and dates, not exclusive use of the name or validity of the framework.

\section{Multimodal Signal Vector \texorpdfstring{$u(t)$}{u(t)}}

Let an audiovisual work of duration $T$ seconds be partitioned into $N$ discrete windows $t \in \{0, 1, \dots, N-1\}$ of duration $\Delta t$. At each window $t$, the work emits a multimodal signal vector:
\begin{equation}
  u(t) = \bigl(u_{\mathrm{ac}}(t),\; u_{\mathrm{vis}}(t),\; u_{\mathrm{ling}}(t)\bigr)
\end{equation}

\subsection{Acoustic Prosody Sub-Vector \texorpdfstring{$u_{\mathrm{ac}}(t)$}{u\_ac(t)}}
Given the downmixed mono $16\text{ kHz}$ PCM16 waveform containing $M$ normalized samples $s_1, \dots, s_M \in [-1, 1]$ within window $t$:
\begin{equation}
  \mathrm{RMS}(t) = \sqrt{\frac{1}{M}\sum_{i=1}^{M} s_i^2}, \quad
  L_{\mathrm{dB}}(t) = 20 \log_{10}\!\bigl(\max(\mathrm{RMS}(t), 10^{-5})\bigr)
\end{equation}
Partitioning window $t$ into $100\text{ ms}$ sub-frames yields the \textbf{Speech Activity Ratio} $\mathrm{VAD}(t) \in [0, 1]$ and intra-window \textbf{Dynamic Range} $\Delta L_{90\text{-}10}(t) = P_{90}(L_{\mathrm{dB}}) - P_{10}(L_{\mathrm{dB}})$. The normalized acoustic intensity $I_{\mathrm{ac}}(t) \in [0, 1]$ is:
\begin{equation}
  I_{\mathrm{ac}}(t) = \mathrm{clip}\!\left(0.55\,\frac{L_{\mathrm{dB}}(t) - L_{\min}}{L_{\max} - L_{\min}} + 0.30\,\mathrm{VAD}(t) + 0.15\,\frac{\Delta L_{90\text{-}10}(t)}{18.0},\; 0,\; 1\right)
\end{equation}

\subsection{Visual Dynamics Sub-Vector \texorpdfstring{$u_{\mathrm{vis}}(t)$}{u\_vis(t)}}
For each sampled frame resized to a $64\times 64$ area-interpolated RGB24 raster ($P = 4096$ pixels), we compute ITU-R BT.709 luminance $Y_p = 0.2126 R_p + 0.7152 G_p + 0.0722 B_p$, spatial contrast $\sigma_Y(t)$, HSV saturation $S(t)$, and the 512-bin quantized color distribution $P_t(b)$.

To measure visual pacing and camera transitions deterministically, we compute the \textbf{Jensen-Shannon Divergence} \cite{Lin1991} (in bits, bounded in $[0, 1]$) between consecutive frames:
\begin{equation}
  D_{\mathrm{JS}}(P_t \parallel P_{t-1}) = \frac{1}{2} D_{\mathrm{KL}}(P_t \parallel M) + \frac{1}{2} D_{\mathrm{KL}}(P_{t-1} \parallel M), \quad M = \frac{1}{2}(P_t + P_{t-1})
\end{equation}
A scene cut is detected whenever $D_{\mathrm{JS}}(t) \ge \tau_{\mathrm{cut}} = 0.30\text{ bits}$. Visual salience $\psi_{\mathrm{vis}}(t) \in [0, 1]$ combines shot flux, spatial contrast, and chromatic entropy $H_{\mathrm{vis}}(t) = -\sum_b P_t(b)\log_2 P_t(b)$:
\begin{equation}
  \psi_{\mathrm{vis}}(t) = \mathrm{clip}\!\left(0.40\,D_{\mathrm{JS}}(t) + 0.35\,\frac{\sigma_Y(t)}{0.32} + 0.25\,\frac{H_{\mathrm{vis}}(t)}{6.5},\; 0,\; 1\right)
\end{equation}

\subsection{Linguistic \& Normative Sub-Vector \texorpdfstring{$u_{\mathrm{ling}}(t)$}{u\_ling(t)}}
From time-aligned speech segments overlapping window $t$, we extract speech cadence $\mathrm{WPM}(t)$, lexical Type-Token Ratio $\mathrm{TTR}(t)$, normative marker density $D_{\mathrm{norm}}(t)$, dialectical connector density $D_{\mathrm{dial}}(t)$, institutional anchoring $A_{\mathrm{inst}}(t) \in [0, 1]$, and active stance vector $z(t) \in [-1, 1]^K$.

\section{Formal Definition of the Six Core NSE Metrics}

All weights and normalization constants below (e.g.\ $0.34$, $0.22$, $0.29$, $0.15$ in $E(t)$; the divisors $0.78$, $0.11$, $0.32$, $6.5$, $18.0$; the damping factor $0.65$) were chosen by the author and have not been calibrated against human judgments or audience data. The composite indices are working definitions, not validated measures.

\subsection{Narrative DNA (\texorpdfstring{$\mathcal{G}_{\mathrm{DNA}}$}{G\_DNA})}
The Narrative DNA is defined as a signed, evidence-linked tripartite graph $\mathcal{G}_{\mathrm{DNA}} = (\mathcal{S} \cup \mathcal{A} \cup \mathcal{B}, \mathcal{E}, W)$ where $\mathcal{S}$ is the set of speaker agents, $\mathcal{A}$ is the set of symbolic archetypes, $\mathcal{B} = \{B_1, \dots, B_K\}$ is the set of core belief propositions, and $W: \mathcal{E} \to [-1, +1]$ assigns signed weights backed by verbatim timestamped quotes.

\subsection{Narrative Energy \texorpdfstring{$E(t) \in [0, 100]$}{E(t)}}
To eliminate zero-collapse while capturing multimodal synergy, Narrative Energy is formulated as a linear-geometric coupling functional:
\begin{align}
  \Phi_{\mathrm{sync}}(t) &= \sqrt{I_{\mathrm{ac}}(t) \cdot I_{\mathrm{ling}}(t)} \\
  E(t) &= 100 \cdot \mathrm{clip}\!\Bigl(0.34\,I_{\mathrm{ac}}(t) + 0.22\,\psi_{\mathrm{vis}}(t) + 0.29\,I_{\mathrm{ling}}(t) + 0.15\,\Phi_{\mathrm{sync}}(t),\; 0,\; 1\Bigr)
\end{align}

\subsection{Narrative Entropy \texorpdfstring{$H_{\mathrm{narr}}(t)$}{H\_narr(t)} (in Shannon Bits)}
Across $K$ belief dimensions, each possessing a positive ($+$) and negative ($-$) pole, let $m_{k,\pm}(t)$ denote the cumulative energy-weighted stance exposure mass up to window $t$, and $p_{k,\pm}(t) = m_{k,\pm}(t) / \sum_{j,s} m_{j,s}(t)$. The Shannon Narrative Stance Entropy \cite{Shannon1948} is:
\begin{equation}
  H_{\mathrm{narr}}(t) = -\sum_{k=1}^{K} \sum_{s \in \{+,-\}} p_{k,s}(t) \log_2 p_{k,s}(t) \in [0, \log_2(2K)]\text{ bits}
\end{equation}
For $K=4$ ($2K = 8$ poles), $H_{\max} = \log_2(8) = 3.00\text{ bits}$, and $\hat{H}_{\mathrm{narr}}(t) = H_{\mathrm{narr}}(t)/3.00 \in [0, 1]$.

\subsection{Narrative Complexity \texorpdfstring{$C(t)$}{C(t)} and Inertia \texorpdfstring{$\mathcal{I}(t)$}{I(t)}}
Complexity $C(t) \in [0, 100]$ and Inertia $\mathcal{I}(t) \in [0, 1]$ are given by:
\begin{align}
  C(t) &= 100 \cdot \mathrm{clip}\!\left(0.40\,\frac{\mathrm{TTR}(t)}{0.78} + 0.35\,\frac{D_{\mathrm{dial}}(t)}{0.11} + 0.25\,\frac{\max(0, |\{k : |z_k(t)| \ge 0.20\}| - 1)}{3.0},\; 0,\; 1\right) \\
  \mathcal{I}(t) &= \mathrm{clip}\!\left(0.55\,A_{\mathrm{inst}}(t) + 0.25\,\frac{1}{K}\sum_{k=1}^{K} |x_k(t)| + 0.20\,\frac{t+1}{N},\; 0,\; 1\right)
\end{align}

\subsection{State-Space Control and Controllability Gramian \texorpdfstring{$W_c$}{W\_c}}
Following Friedkin-Johnsen opinion dynamics \cite{Friedkin1990}, we \emph{hypothesize} that a collective belief state $x(t) \in [-1, 1]^K$ evolves as a discrete-time Linear Time-Varying (LTV) system driven by stance input $z(t)$. No audience belief state is measured in this work; $x(t)$ is a model variable:
\begin{equation}
  x_k(t+1) = \bigl(1 - \gamma_k(t)\bigr) x_k(t) + \gamma_k(t) z_k(t), \qquad
  \gamma_k(t) = \eta \cdot \frac{E(t)}{100} \cdot \bigl(1 - 0.65\,\mathcal{I}_k(t)\bigr)
\end{equation}
with $\eta = 0.18$. Over horizon $[0, N]$, let $\Phi_k(N, t+1) = \prod_{\tau=t+1}^{N-1} (1 - \gamma_k(\tau))$ and $\tilde{b}_k(t) = \gamma_k(t)(0.35 + 0.65|z_k(t)|)$. The diagonal \textbf{Controllability Gramian} is:
\begin{equation}
  W_{c,kk}(0, N) = \sum_{t=0}^{N-1} \Bigl(\Phi_k(N, t+1)\,\tilde{b}_k(t)\Bigr)^2, \qquad
  \kappa(W_c) = \frac{\lambda_{\max}(W_c)}{\lambda_{\min}(W_c)}
\end{equation}

\begin{table}[!htbp]
\centering
\caption{Results Across Speaker Turns in \texttt{CORPUS-001} (\emph{Insumisi\'on Debate}, Spain, 1994, $194.84\text{ s}$, $N=39$).}
\label{tab:exp001}
\small
\resizebox{\textwidth}{!}{\begin{tabular}{llccccccc}
\toprule
\textbf{Turn ID} & \textbf{Speaker / Visual Anchor} & \textbf{Span (s)} & \textbf{dBFS} & \textbf{WPM} & \textbf{Mean $E$} & \textbf{Comp. $C$} & \textbf{Inertia $\mathcal{I}$} & \textbf{Stance $(B_1, B_2, B_3, B_4)$} \\
\midrule
\texttt{ENV\_POSTER} & Activist Poster Close-up & $0\text{--}5$ & $-26.36$ & $312.0^*$ & $52.40$ & $65.00$ & $0.088$ & $(+0.90, -0.85, -0.20, +0.80)$ \\
\texttt{SPK\_01} & Student 1 (Metallica shirt) & $5\text{--}45$ & $\mathbf{-16.96}$ & $\mathbf{203.7}$ & $\mathbf{72.90}$ & $75.67$ & $0.151$ & $(+0.93, -0.66, -0.38, +0.67)$ \\
\texttt{SPK\_02} & Student 2 (White/pink shirt) & $45\text{--}90$ & $-19.83$ & $141.8$ & $63.86$ & $\mathbf{87.90}$ & $0.539$ & $(-0.48, +0.77, +0.91, -0.32)$ \\
\texttt{SPK\_03} & Student 3 (World map wall) & $90\text{--}110$ & $-21.34$ & $201.0$ & $66.45$ & $76.15$ & $\mathbf{0.668}$ & $(-0.55, +0.95, +0.70, -0.60)$ \\
\texttt{SPK\_04} & Student 4 (Navy NY shirt) & $110\text{--}130$ & $-18.64$ & $192.0$ & $65.08$ & $59.47$ & $0.536$ & $(+0.15, +0.60, +0.20, +0.85)$ \\
\texttt{SPK\_05} & Student 5 (Plaid shirt, hall) & $130\text{--}194.8$ & $-22.21$ & $185.1$ & $57.18$ & $68.39$ & $0.663$ & $(-0.25, +0.72, +0.89, -0.04)$ \\
\bottomrule
\end{tabular}}
\end{table}

\section{Experiment \texttt{EXP-001}: Descriptive Application}

We first evaluated the framework on \texttt{06-FCE-V-C110-007-C.mp4} ($194.84\text{ s}$, $N=39$ windows), an archival Spanish university debate on compulsory military conscription (\emph{la mili}) vs.\ civil disobedience (\emph{insumisi\'on}). Table~\ref{tab:exp001} summarizes the multimodal telemetry across the six turns over four belief dimensions, with stance vectors coded by a single coder (no agreement statistic yet): $B_1$ (\emph{Anti-Militarist Abolitionism}), $B_2$ (\emph{Civic \& Legal Compliance}), $B_3$ (\emph{Pragmatic Social Service}), and $B_4$ (\emph{Individual Conscience Autonomy}).

Observations from \texttt{EXP-001}:
\begin{enumerate}
  \item \textbf{Visual Transitions:} with one frame every 5 s, $D_{\mathrm{JS}} \ge 0.30\text{ bits}$ flagged Windows \#1, \#18, and \#26, which coincide with setting changes seen on inspection (Window \#0 is marked by convention). No annotated shot-boundary reference exists, so error rates are unknown.
  \item \textbf{Energy vs.\ Inertia Asymmetry:} \texttt{SPK\_01} achieved the highest Narrative Energy ($E = 76.54$ at Window \#7; mean $72.90$) due to loud prosody ($-16.96\text{ dBFS}$) and rapid cadence ($203.7\text{ WPM}$), but the lowest Inertia ($\mathcal{I} = 0.151$), while \texttt{SPK\_03}, $4.4\text{ dB}$ softer, received the highest ($\mathcal{I} = 0.668$). Inertia is dominated by coded institutional anchoring (weight $0.55$), so this contrast follows from the coding, not from an independent measurement.
\end{enumerate}

\begin{table}[!htbp]
\centering
\caption{Three-Work Comparison (\texttt{EXP-002}) Across $513.6\text{ s}$ ($77$ Windows).}
\label{tab:exp002}
\small
\resizebox{\textwidth}{!}{\begin{tabular}{lccc}
\toprule
\textbf{Metric / Dimension} & \textbf{\texttt{CORPUS-001}: Insumisi\'on Debate (ES)} & \textbf{\texttt{CORPUS-002}: LBJ ``Daisy'' (EN)} & \textbf{\texttt{CORPUS-003}: Reagan Challenger (EN)} \\
\midrule
Communicative Regime & Polyphonic Civic Debate & Acute Fear-Appeal Spot & Institutional Crisis Consolation \\
Duration \& Windows & $194.84\text{ s}$ ($39\times 5\text{ s}$) & $59.83\text{ s}$ ($12\times 5\text{ s}$) & $258.93\text{ s}$ ($26\times 10\text{ s}$) \\
Mean Narrative Energy $\bar{E}$ & $\mathbf{63.58}$ & $50.51$ & $56.45$ \\
Peak Narrative Energy $E_{\max}$ & $76.54$ (Win \#7) & $\mathbf{83.29}$ (Win \#6) & $61.98$ (Win \#10) \\
Energy Dynamic Range ($E_{\max}-E_{\min}$) & $33.91$ & $\mathbf{59.40}$ ($23.9\text{--}83.3$) & $31.76$ \\
Final Stance Entropy $H(T)$ & $\mathbf{2.7574\text{ bits}}$ ($91.9\%$) & $2.2292\text{ bits}$ ($74.3\%$) & $\mathbf{2.0240\text{ bits}}$ ($67.5\%$) \\
Mean Narrative Inertia $\bar{\mathcal{I}}$ & $0.5019$ & $0.4420$ & $\mathbf{0.7266}$ (peaks at $\mathbf{0.8921}$) \\
Flagged transitions (excl.\ Win \#0) & $3$ ($\max 0.490$) & $2$ ($\mathbf{\max 0.577}$) & $\mathbf{0}$ ($\max 0.042$) \\
Final Polarization Index $P(T)$ & $\mathbf{0.5515}$ & $0.4200$ & $\mathbf{0.0617}$ \\
Control Condition $\kappa(W_c)$ & $2.738$ & $1.783$ & $1.457$ \\
\bottomrule
\end{tabular}}
\end{table}

\section{Experiment \texttt{EXP-002}: Three-Work Comparison}

To see how the same definitions behave across languages and purposes, we applied them without recalibration to two public-domain U.S.\ broadcasts alongside \texttt{CORPUS-001} (Table~\ref{tab:exp002}):
\begin{itemize}
  \item \textbf{\texttt{CORPUS-002} (\emph{LBJ ``Daisy'' Nuclear Spot}, 1964, $59.83\text{ s}$):} Windows \#0--\#3 establish a quiet meadow baseline ($E \in [23.89, 32.82]$, $-28.2\text{ dBFS}$). At Window \#6 ($30\text{--}35\text{ s}$, nuclear detonation), visual divergence spikes to $D_{\mathrm{JS}} = 0.5766\text{ bits}$, acoustic level hits $-13.68\text{ dBFS}$, and Narrative Energy rises $3.49\times$ to $\mathbf{83.29/100}$.
  \item \textbf{\texttt{CORPUS-003} (\emph{Reagan Challenger Address}, 1986, $258.93\text{ s}$):} Shows no flagged transitions ($\max D_{\mathrm{JS}} = 0.0419\text{ bits}$), lowest stance entropy ($2.0240\text{ bits}$), near-zero polarization ($0.0617$), and maximum Institutional Inertia ($\mathcal{I}(T) = 0.8921$).
\end{itemize}
In the normalized 5D NSE Metric Space $\mathbf{m} = (\bar{E}/100, E_{\max}/100, \hat{H}(T), \bar{\mathcal{I}}, P(T))$, the pairwise Euclidean distances are $\|\mathbf{m}_{\mathrm{C1}} - \mathbf{m}_{\mathrm{C2}}\|_2 = 0.2711$, $\|\mathbf{m}_{\mathrm{C2}} - \mathbf{m}_{\mathrm{C3}}\|_2 = 0.5128$, and $\|\mathbf{m}_{\mathrm{C1}} - \mathbf{m}_{\mathrm{C3}}\|_2 = \mathbf{0.6135}$ (simulated terminal state distance $1.0045$). With one work per purpose, these distances describe three works and cannot show genre-level separation.

\begin{table}[!htbp]
\centering
\caption{Simulated Terminal States for Selected Block Orders Among All $6! = 720$ Permutations in \texttt{EXP-003} (\texttt{CORPUS-001}). Model simulations; no viewers involved.}
\label{tab:exp003}
\small
\resizebox{\textwidth}{!}{\begin{tabular}{llcccccc}
\toprule
\textbf{Control Objective} & \textbf{Block Order $\pi^*$} & $x_1(T)$ & $x_2(T)$ & $x_3(T)$ & $x_4(T)$ & \textbf{Norm $\|x(T)\|_2$} & \textbf{Shift $\|\Delta x\|_2$} \\
\midrule
Original Broadcast & \texttt{P $\to$ S1 $\to$ S2 $\to$ S3 $\to$ S4 $\to$ S5} & $-0.1445$ & $+0.5524$ & $+0.6446$ & $+0.0391$ & $0.8620$ & $0.0000$ \\
$\pi^*_{\mathrm{anti}}$ (Max Abolitionism) & \texttt{S3 $\to$ S2 $\to$ S5 $\to$ S4 $\to$ P $\to$ S1} & $\mathbf{+0.5130}$ & $-0.1717$ & $-0.0304$ & $+0.4568$ & $0.7087$ & $\mathbf{1.2597}$ \\
$\pi^*_{\mathrm{legal}}$ (Max Legal Compliance) & \texttt{P $\to$ S1 $\to$ S4 $\to$ S5 $\to$ S2 $\to$ S3} & $-0.2601$ & $\mathbf{+0.5830}$ & $+0.6705$ & $-0.1552$ & $0.9387$ & $0.2296$ \\
$\pi^*_{\mathrm{eq}}$ (Minimum Norm) & \texttt{S4 $\to$ P $\to$ S5 $\to$ S2 $\to$ S1 $\to$ S3} & $+0.2262$ & $+0.0606$ & $+0.1787$ & $+0.1430$ & $\mathbf{0.3274}$ & $0.7792$ \\
$\pi^*_{\mathrm{aut}}$ (Max Autonomy) & \texttt{S3 $\to$ S2 $\to$ S5 $\to$ S1 $\to$ P $\to$ S4} & $+0.3958$ & $-0.0083$ & $+0.0669$ & $\mathbf{+0.4930}$ & $0.6358$ & $1.0705$ \\
\bottomrule
\end{tabular}}
\end{table}

\section{Experiment \texttt{EXP-003}: Counterfactual Simulations}

All results in this section are simulations of the model of Section~3 under modified inputs.
\begin{enumerate}
  \item \textbf{All $6! = 720$ block orders (\texttt{CORPUS-001}):} reordering the six speaker blocks moves every simulated belief dimension across zero, e.g.\ $x_1(T) \in [-0.2603, +0.5130]$ and $x_2(T) \in [-0.1795, +0.5830]$. The minimum-norm order $\pi^*_{\mathrm{eq}}$ (Table~\ref{tab:exp003}) gives $\|x^*(T)\|_2 = 0.3274$ vs.\ $0.8620$ for the broadcast order ($-62.0\%$). Reading a smaller norm as ``less bias'' is a modeling choice; the result shows the model's order sensitivity, not an effect on viewers.
  \item \textbf{Attenuating the non-verbal peak (\texttt{CORPUS-002}):} attenuating Windows \#4--\#8 of \emph{``Daisy''} with the transcript unchanged lowers the simulated $x_2(T)$ from $+0.2624$ to $+0.1841$; within the model, $29.84\%$ of the displacement is attributable to those acoustic and visual inputs. The share depends on the uncalibrated channel weights.
  \item \textbf{Removing institutional anchoring (\texttt{CORPUS-003}):} setting coded anchoring to zero lowers terminal inertia by $0.4657$ ($0.8921 \to 0.4264$). Injecting a $30\text{ s}$ high-energy critique at $t = 30$, $120$, or $210\text{ s}$ lowers the simulated $x_2(T)$ from $+0.667$ to $+0.5013$, $+0.3836$, and $+0.2321$; later injections have larger simulated effects, reflecting the recency of the update rule.
\end{enumerate}

\section{Experiment \texttt{EXP-004}: Estimator Self-Consistency}

No measured belief trajectories exist, so EXP-004 generates target trajectories with the model itself using fixed parameters $\theta^* = (0.34, 0.22, 0.29, 0.15;\ \beta_0 = 0.18)$, adds Gaussian noise, and estimates $\theta$ by ridge regression. The estimator recovers $\theta^*$ exactly without noise and with rollout $R^2 = 0.9999$ at $25\text{ dB}$ SNR; leave-one-corpus-out fits on the simulated trajectories reach $R^2_{\mathrm{OOS}} \ge 0.9997$. Because every corpus is generated by the same model and parameters, these values verify the estimator and identifiability, not the model or its generalization.

\section{Limitations and Next Steps}

\textbf{Limitations.} (1) The belief state $x(t)$ is never observed; all statements about belief change, control, or montage effects are properties of the model. (2) Stance and anchoring labels were coded by one coder. (3) Composite weights are uncalibrated. (4) Three works, one per purpose. (5) One frame per window; no annotated shot-boundary reference. (6) EXP-004 uses synthetic trajectories.

\textbf{Next steps.} Two independent coders with a preregistered agreement threshold; annotated shot boundaries to measure detector precision and recall; a larger corpus with several works per purpose and a held-out comparison; a preregistered audience study comparing montage orders of \texttt{CORPUS-001}; and calibration of composite weights against human ratings.

\section{Conclusion}

The measurement layer of NSE turns audiovisual works into reproducible, time-aligned profiles of acoustic, visual, and linguistic signals, and the three-work application shows that these profiles capture visible differences in pacing and delivery. The state-space layer is a formal, executable hypothesis about how such signals might shift audience beliefs; it has not been tested against audiences.

\paragraph{Revision note (v1.1).} v1.0 described the results as validation of a discipline, as proof that montage order controls belief, and as cross-domain generalization. v1.1 withdraws those claims and corrects the visual transition counts, which in v1.0 included Window \#0 by convention. Numerical outputs are otherwise unchanged.

\bibliographystyle{plain}
\bibliography{references}

\end{document}
